% ==============================================================
% Kapitel 05: Implementation
% ==============================================================

\section{Implementation}

\subsection{Server}

\subsubsection{Zugriffsberechtigungen}
\label{sec:impl_zugriffsberechtigungen}

Die Zugriffskontrolle im Backend basiert auf FastAPIs 
Dependency-Injection-System. Rollen werden einmal definiert und dann 
über die gesamte Anwendung verwendet. Die Hierarchie ist so aufgebaut, 
dass höhere Rollen automatisch die Berechtigungen der darunterliegenden Rollen einschließen.

\begin{CodeBox}[language=Python]
ROLE_HIERARCHY = {
    UserRole.VIEWER: {
        UserRole.VIEWER,
        UserRole.TEAM_LEAD,
        UserRole.ADMIN,
    },
    UserRole.TEAM_LEAD: {
        UserRole.TEAM_LEAD,
        UserRole.ADMIN,
    },
    UserRole.ADMIN: {
        UserRole.ADMIN,
    },
}

def require_role(required_role: UserRole):
    def role_checker(user: CurrentUser):
        user_roles = {UserRole(role) for role in user["roles"]}
        allowed_roles = ROLE_HIERARCHY[required_role]
        if user_roles.isdisjoint(allowed_roles):
            raise exception.InsufficientPermissions
        return user
    return role_checker

AdminUser = Annotated[dict, Depends(require_role(UserRole.ADMIN))]
TeamLeadUser = Annotated[dict, Depends(require_role(UserRole.TEAM_LEAD))]
ViewerUser = Annotated[dict, Depends(require_role(UserRole.VIEWER))]
\end{CodeBox}

Die oben definierten Typen können direkt als Parameter in Endpunkt-Funktionen verwendet werden.
FastAPI löst die Dependency automatisch auf und führt die Rollenprüfung vor der eigentlichen Endpunktlogik aus.

\begin{CodeBox}[language=Python, caption={Verwendung der Rollendependencies in Endpunkten}]
@router.get("/internal/verify/teamlead")
def verify_teamlead(request: Request, response: Response,
                    current_user: TeamLeadUser):
    user_resp = requests.get(
        f"{USER_URL}/api/users/me",
        headers={"Authorization": request.headers.get("Authorization")}
    )
    user_data = user_resp.json()
    response.headers["X-User-Id"] = current_user["sub"]
    response.headers["X-Team-Id"] = str(user_data.get("team_id"))
    response.headers["X-Role"] = current_user["roles"][0]
    return {"active": True, **current_user}
\end{CodeBox}

Der \texttt{verify\_teamlead}-Endpunkt hat eine Besonderheit.
Er sendet die Team-ID des Benutzers als HTTP-Antwort-Header zurück, 
damit Nginx diese Informationen an die Ziel-Services weiterreichen kann.

\subsubsection{API-Gateway}
\label{sec:impl_nginx}

Das Nginx \gls{API}-Gateway wird über eine zentrale Konfigurationsdatei 
definiert. Die grundlegende Struktur besteht aus der Definition 
der Upstream-Services sowie einem Server-Block, der alle 
Routing- und Authentifizierungsregeln enthält.

\paragraph{Upstream-Konfiguration}

Jeder Microservice wird als eigener Upstream registriert, 
damit Nginx die Anfragen richtig weiterleiten kann. Hier der User Service als Beispiel:
\begin{CodeBox}[language=bash, caption={Upstream-Definition der Microservices}]
upstream user_service {server user-service:8001;}
\end{CodeBox}

\newpage

\paragraph{Interne Authentifizierungsendpunkte}
Die Authentifizierung wird über drei interne Endpunkte realisiert, 
die ausschließlich vom \texttt{auth\_request}-Modul aufgerufen 
werden und von außen nicht erreichbar sind. Der einzige Unterschied zwischen den drei 
ist die URL, die auf den Auth Service zeigt, welcher die jeweilige Rolle überprüft.
Hier der Endpunkt für die Admin-Authentifizierung als Beispiel:
\begin{CodeBox}[language=bash]
location = /auth_admin {
    internal;
    proxy_pass http://auth_service/api/auth/internal/verify/admin;
    proxy_pass_request_body off;
    proxy_set_header Content-Length "";
}
\end{CodeBox}

\paragraph{Routing mit Authentifizierung}
Jeder Endpunkt verwendet \texttt{auth\_request}, um die Berechtigung zu prüfen. Das folgende 
Beispiel zeigt das Routing für den User Service, der ausschließlich 
Administratoren zugänglich ist:

\begin{CodeBox}[language=bash]
location ~ ^/users(/.*)?$ {
    auth_request /auth_admin;
    proxy_pass http://user_service/api/users$1$is_args$args;
}
\end{CodeBox}

\paragraph{API-Key Bypass}
Für die geplante Anbindung an die Zeitnehmer-App kann die Authentifizierung mittels \gls{API}-Key 
umgangen werden:

\begin{CodeBox}[language=bash]
map $http_x_api_key $skip_auth {
    default 0;
    "this-is-not-the-key" 1;
}
location = /auth_admin {
    internal;
    if ($skip_auth = 1) { return 204; }
    proxy_pass http://auth_service/api/auth/internal/verify/admin;
}
\end{CodeBox}

\subsubsection{CRUD-Operationen}
\label{sec:impl_crud}
Die \gls{CRUD}-Operationen sind in allen Services nach demselben Muster 
aufgebaut: Ein Router-Modul definiert die HTTP-Endpunkte und 
delegiert die eigentliche Logik an ein separates \gls{CRUD}-Modul. 
Dieses Muster wird am Beispiel des Team Service erläutert.

\paragraph{Router}
Der Router definiert die Endpunkte und deren HTTP-Methoden. 
Die Datenbankverbindung wird über FastAPIs Dependency-Injection 
als \texttt{SessionDep} bereitgestellt, sie wird automatisch 
geöffnet, an den Endpunkt übergeben und nach dem Request 
wieder geschlossen. Auf der HTTP-Ebene ändert sich dann immer nur die Methode,
die der Router aufruft (POST, GET, PUT, DELETE).
Hier der Create-Endpunkt des Team Service als Beispiel:

\begin{CodeBox}[language=Python]
@router.post("/", response_model=TeamResponse)
def create_team(db: SessionDep, team: TeamCreate):
    return crud.create_team(db=db, team=team)
\end{CodeBox}

\paragraph{Erstellen}
Beim Erstellen eines Teams wird das eingehende Pydantic-Schema 
direkt in ein SQLAlchemy-Modell umgewandelt und in der Datenbank 
gespeichert. Bei einem Fehler wird die Transaktion zurückgerollt:
\begin{CodeBox}[language=Python, caption={Erstellen eines Teams}]
def create_team(*, db: SessionDep, team: TeamCreate):
    try:
        db_team = Team(**team.model_dump(exclude_unset=True))
        db.add(db_team)
        db.commit()
        db.refresh(db_team)
        return db_team
    except Exception as exc:
        db.rollback()
        raise ServiceError() from exc
\end{CodeBox}

\newpage

\paragraph{Aktualisieren}

Beim Aktualisieren werden nur die Felder überschrieben, die 
tatsächlich in der Anfrage enthalten sind. Das 
\texttt{exclude\_unset=True}-Flag von Pydantic stellt sicher, 
dass nicht gesetzte Felder nicht auf ihre Standardwerte 
zurückgesetzt werden:

\begin{CodeBox}[language=Python, caption={Aktualisieren eines Teams}]
def update_team(*, db: SessionDep, team_id: int,
                team_update: TeamUpdate, request: Request):
    check_team_permissions(db=db, team_id=team_id, request=request)
    try:
        db_team = get_team(db=db, team_id=team_id, request=request)
        update_data = team_update.model_dump(
            exclude_unset=True,
            exclude={"id"},
        )
        for field, value in update_data.items():
            setattr(db_team, field, value)
        db.commit()
        db.refresh(db_team)
        return db_team
    except EntityDoesNotExistError:
        db.rollback()
        raise
    except Exception as exc:
        db.rollback()
        raise ServiceError() from exc
\end{CodeBox}

\newpage

\paragraph{Löschen}

Beim Löschen wird das angeforderte Team aus der Datenbank entfernt.
Zuvor wird eine Berechtigungsprüfung durchgeführt, um sicherzustellen, dass der Zugriff erlaubt ist.

\begin{CodeBox}[language=Python, caption={Löschen eines Teams mit Vorabprüfung}]
def delete_team(*, db: SessionDep, team_id: int, request: Request):
    check_team_permissions(db=db, team_id=team_id, request=request)
    try:
        db_team = get_team(db=db, team_id=team_id, request=request)
        db.delete(db_team)
        db.commit()
        return db_team
    except EntityDoesNotExistError:
        db.rollback()
        raise
    except Exception as exc:
        db.rollback()
        raise ServiceError() from exc
\end{CodeBox}

\paragraph{Rollenbasierte Berechtigungsprüfung}

Zugriffe werden zusätzlich durch eine 
serviceinterne Berechtigungsprüfung abgesichert. Teamleiter 
dürfen ausschließlich auf ihr eigenes Team zugreifen,
die Team-ID des anfragenden Benutzers wird dabei aus dem 
vom Nginx API-Gateway weitergeleiteten \texttt{X-Team-Id}-Header 
ausgelesen:

\begin{CodeBox}[language=Python, caption={Rollenbasierte Berechtigungsprüfung im Team Service}]
def check_team_permissions(*, db: SessionDep,
                           team_id: int | None = None,
                           request: Request):
    role = request.headers.get("X-Role")
    user_team_id = request.headers.get("X-Team-Id")
    if role == "TEAM_LEAD" and team_id is not None:
        if int(team_id) != int(user_team_id):
            raise InsufficientPermissions(
                "Teamleads can only operate on their own team."
            )
\end{CodeBox}

\subsubsection{Authentifizierung mit Keycloak}
\label{sec:impl_keycloak}
Die Authentifizierung funktioniert mit JWT-Tokens, die von Keycloak 
ausgestellt und vom Auth Service validiert werden. Keycloak 
übernimmt dabei die vollständige Verwaltung von Benutzerkonten, 
Passwörtern und Rollen. Der Auth Service fungiert als 
Vermittler zwischen Keycloak und dem restlichen System.

\paragraph{Login}

Beim Login sendet die Webanwendung Benutzername und Passwort 
an den Auth Service, der diese über den OAuth2 Resource Owner 
Password Flow an Keycloak weiterleitet. Bei Erfolg liefert 
Keycloak ein Access- sowie ein Refresh-Token zurück, die 
in der Webanwendung im \texttt{localStorage} gespeichert werden.

\begin{figure}[h]
    \centering
    \includegraphics[width=\textwidth]{login_ablauf.png}
    \caption{Sequenzdiagramm: Login-Ablauf}
    \label{fig:seq_login}
\end{figure}

\newpage

\paragraph{Authentifizierte Anfrage}

Bei jeder API-Anfrage prüft Nginx über das 
\texttt{auth\_request}-Modul, ob der mitgesendete JWT gültig 
ist. Der Auth Service validiert den Token lokal anhand der 
gecachten öffentlichen Schlüssel von Keycloak, Keycloak 
selbst wird dabei nicht eingebunden. Nach erfolgreicher 
Validierung wird die Anfrage mit den Benutzerinformationen 
als Header an den Ziel-Service weitergeleitet.

\begin{figure}[h]
    \centering
    \includegraphics[width=\textwidth]{AuthentifizierteAnfrage.png}
    \caption{Sequenzdiagramm: Authentifizierte Anfrage}
    \label{fig:seq_request}
\end{figure}

\newpage

\paragraph{Token-Refresh}

Stellt die Webanwendung fest, dass das Access-Token abgelaufen 
ist, wird vor der eigentlichen Anfrage automatisch ein 
Token-Refresh durchgeführt. Ist auch das Refresh-Token 
abgelaufen, wird die Sitzung beendet und der Benutzer zur 
Login-Seite weitergeleitet.

\begin{figure}[h]
    \centering
    \includegraphics[width=\textwidth]{token_refresh.png}
    \caption{Sequenzdiagramm: Token-Refresh}
    \label{fig:seq_refresh}
\end{figure}

\subsubsection{Inter-Service-Kommunikation}
\label{sec:impl_interservice}

Die Services kommunizieren untereinander über synchrone 
HTTP-Aufrufe mittels der Python-\texttt{requests}-Bibliothek\cite{requests}. 
Die URLs der anderen Services werden dabei über Umgebungsvariablen 
konfiguriert und beim Start des Containers eingelesen, sodass 
keine URLs hart im Code hinterlegt sind.

\paragraph{Validierung von Fremdbeziehungen}
Wie in~\ref{sec:server} beschrieben, existieren keine echten Datenbankfremdschlüssel über Servicegrenzen hinweg.
Stattdessen werden Referenzen auf andere Services vor dem Speichern explizit validiert. Der Attempt Service prüft 
beispielsweise vor dem Erstellen eines Versuchs, ob das 
referenzierte Team, der Fahrer und der Bewerb tatsächlich 
existieren:

\begin{CodeBox}[language=Python, caption={Validierung von Fremdbeziehungen vor dem Speichern}]
def _validate_team(team_id: int):
    resp = requests.get(f"{TEAM_URL}/api/teams/{team_id}")
    if resp.status_code == 404:
        raise EntityDoesNotExistError(f"Team {team_id} does not exist")
    if resp.status_code in (401, 403):
        raise AuthenticationFailed(f"Unauthorized to fetch team {team_id}")
    if resp.status_code != 200:
        raise ServiceError(f"Failed to fetch team {team_id}: {resp.text}")

def create_attempt(*, db: SessionDep, attempt: AttemptCreate):
    _validate_team(attempt.team_id)
    _validate_driver(attempt.driver_id)
    _validate_challenge(attempt.challenge_id)
    # ...
\end{CodeBox}

\paragraph{Auslösen von Folgeoperationen}

Nach dem Speichern eines Versuchs löst der Attempt Service 
automatisch die Punkteberechnung im Score Service aus. 
Werden dabei optional auch Strafpunkte mitgeliefert, werden 
diese zuerst angelegt:

\begin{CodeBox}[language=Python]
if attempt.penalty_count and attempt.penalty_type:
    penalty_payload = {
        "attempt_id":    db_attempt.id,
        "count":         attempt.penalty_count,
        "penalty_type_id": attempt.penalty_type,
    }
    pen_resp = requests.post(
        f"{SCORE_URL}/api/penalties/", json=penalty_payload

score_resp = requests.post(
    f"{SCORE_URL}/api/scores/",
    json={"attempt_id": db_attempt.id}
)
\end{CodeBox}

\subsection{Web App}
\label{sec:impl_webapp}

\subsubsection{Kommunikation mit dem Server}
\label{sec:impl_webapp_api}

Die Kommunikation mit dem Backend erfolgt über typisierte 
API-Client-Module, die für jeden Service ein eigenes Objekt 
bereitstellen. Alle authentifizierten Anfragen laufen über 
die zentrale \texttt{authenticatedFetch}-Funktion, die den 
Token automatisch als \texttt{Authorization}-Header 
beifügt:

\begin{CodeBox}[, caption={Zentrale Fetch-Funktion mit automatischer Authentifizierung}]
export async function authenticatedFetch(
  url: string,
  options: RequestInit = {}
): Promise<Response> {
  const token = await AuthService.getValidToken()

  if (!token) {
    window.dispatchEvent(new Event('auth:session-expired'))
    throw new Error('Not authenticated')
  }

  const headers = new Headers(options.headers)
  headers.set('Authorization', `Bearer ${token}`)
  if (!(options.body instanceof FormData)) {
    headers.set('Content-Type', 'application/json')
  }
  return fetch(url, { ...options, headers })
}
\end{CodeBox}

\newpage

Die API-Client-Module abstrahieren die HTTP-Aufrufe hinter 
typisierten Funktionen. Das folgende Beispiel zeigt die 
\texttt{createTeam}-Funktion des Team Service:

\begin{CodeBox}[, caption={Typisierter API-Client am Beispiel des Team Service}]
  async createTeam(data: TeamCreate): Promise<Team> {
    const response = await authenticatedFetch(`${API_BASE_URL}/teams/`, {
      method: 'POST',
      body: JSON.stringify(data),
    })
    if (!response.ok) throw new Error('Failed to create team')
    return response.json()
  },
\end{CodeBox}

\subsubsection{Rollenbasierte Zugriffskontrolle}
\label{sec:impl_webapp_permissions}

Die Zugriffskontrolle auf der Clientseite wird über ein 
zentrales Berechtigungsmodul gesteuert. Für jede Rolle 
ist definiert, welche Tabs sichtbar und welche Aktionen 
erlaubt sind.
Hier ein Ausschnitt aus der Berechtigungsdefinition:

\begin{CodeBox}[, caption={Rollenbasierte Berechtigungsdefinition}]
const rolePermissions: Record<UserRole, RolePermissions> = {
  admin: {
    allowedTabs: ['leaderboard', 'attempts', 'teams',
                  'challenges', 'users', 'export', 'login'],
    canViewAttempts: true,
    canEditTeams: true,
    canEditUsers: true,
    canEditChallenges: true,
    canExport: true,
  },
\end{CodeBox}

In der Hauptkomponente wird bei jedem Tab-Wechsel geprüft, 
ob die Rolle des aktuellen Benutzers den Zugriff erlaubt.

\begin{CodeBox}[, caption={Tab-Zugriffskontrolle in der Hauptkomponente}]
const handleTabChange = (tab: Tabs) => {
  if (canAccessTab(currentUser?.role || null, tab)) {
    setActiveTab(tab)
  }
}
\end{CodeBox}

